%  LaTeX support: latex@mdpi.com
%=================================================================
\documentclass[symmetry,article,submit,moreauthors]{Definitions/mdpi}

% MDPI internal commands
\firstpage{1}
\makeatletter
\setcounter{page}{\@firstpage}
\makeatother
\pubvolume{1}
\issuenum{1}
\articlenumber{1}
\pubyear{2026}
\copyrightyear{2026}
\datereceived{}
\daterevised{}
\dateaccepted{}
\datepublished{}

%=================================================================
\Title{Symmetric vs. Asymmetric Knowledge Distillation: A Comparative Study for Image Classification}

% Author ORCID
\newcommand{\orcidauthorA}{0000-0000-0000-000X}
\newcommand{\orcidauthorB}{0000-0000-0000-000X}

% Authors
\Author{Firstname Lastname $^{1}$\orcidA{}, Firstname Lastname $^{2,}$*}

\AuthorNames{Firstname Lastname, Firstname Lastname}

\address{%
$^{1}$ \quad Affiliation 1; e-mail@e-mail.com\\
$^{2}$ \quad Affiliation 2; e-mail@e-mail.com}

\corres{Correspondence: e-mail@e-mail.com}

%=================================================================
\abstract{Knowledge distillation (KD) is a widely adopted technique for compressing deep neural networks. Despite the proliferation of KD methods, a systematic understanding of their underlying mechanisms remains elusive. In this paper, we propose a perspective that categorizes KD methods along a symmetry spectrum: symmetric methods directly align teacher and student representations through one-to-one feature matching, while asymmetric methods preserve relational structures among samples without requiring direct feature correspondence. Through extensive experiments on CIFAR-100 with ResNet-34 as teacher and ResNet-18 as student, we compare six representative KD methods (KD, FitNet, AT, SP, CC, RKD). Our results reveal three key findings: (1) both symmetric and asymmetric methods can effectively improve student performance; (2) symmetric methods converge faster while asymmetric methods achieve competitive or higher peak accuracy; (3) the boundary between symmetric and asymmetric methods is a spectrum rather than a dichotomy, with several methods exhibiting hybrid characteristics. This work provides practitioners with a complementary conceptual framework for understanding and selecting knowledge distillation strategies.}

\keyword{knowledge distillation; symmetric learning; asymmetric learning; model compression; image classification; representation learning}

%=================================================================
\begin{document}

%%%%%%%%%%%%%%%%%%%%%%%%%%%%%%%%%%%%%%%%%%
\section{Introduction}

Deep convolutional neural networks (CNNs) have achieved remarkable success in image classification tasks, with architectures such as ResNet~\cite{ref1}, DenseNet~\cite{ref2}, and EfficientNet~\cite{ref3} pushing the boundaries of accuracy. However, the superior performance of these models comes at the cost of substantial computational and memory requirements, making their deployment on resource-constrained devices challenging~\cite{ref4}. This trade-off has motivated extensive research on model compression techniques, including pruning~\cite{ref5}, quantization~\cite{ref6}, and knowledge distillation (KD)~\cite{ref7}.

Knowledge distillation~\cite{ref7} addresses this challenge by transferring knowledge from a large, pre-trained teacher network to a smaller student network. The student is trained to mimic the teacher's output distribution by minimizing the Kullback-Leibler (KL) divergence between their softened probability distributions. Since this seminal work, numerous extensions have been proposed, which can be broadly categorized along multiple dimensions: output-level vs.\ feature-level~\cite{ref8,ref9}, logit-based vs.\ structure-based~\cite{ref10,ref11}, and instance-level vs.\ relational-level~\cite{ref12}.

In this paper, we propose a complementary perspective that cuts across these existing categorizations: \textbf{the symmetry of the knowledge transfer mechanism}. Specifically, some KD methods perform a direct, one-to-one alignment between teacher and student representations. We term these \textbf{symmetric} methods, as they assume a symmetric correspondence between the representation spaces of the two networks. Examples include the original KD (output distribution matching), FitNet~\cite{ref8} (intermediate feature regression), and Attention Transfer (AT)~\cite{ref9} (attention map alignment). In contrast, other methods preserve the relational structure among samples without requiring direct feature correspondence. We term these \textbf{asymmetric} methods. Examples include Similarity Preserving (SP)~\cite{ref10}, Correlation Congruence (CC)~\cite{ref11}, and Relational Knowledge Distillation (RKD)~\cite{ref12}.

This distinction has practical implications. When teacher and student networks share similar architectures, their intermediate representations are likely to exhibit similar structures, making symmetric alignment natural and effective. However, when the architecture gap is large, direct feature alignment may force the student to learn representations that are not naturally suited to its capacity. Asymmetric methods, by focusing on relational structures, may be more robust to such disparities. We note that this taxonomy represents a spectrum rather than a strict dichotomy---some methods incorporate both symmetric and asymmetric components, as discussed in Section~\ref{sec:discussion}.

Our main contributions are: (1) a new perspective for understanding KD methods based on the symmetry of their knowledge transfer mechanisms; (2) a comprehensive empirical comparison of six representative KD methods (KD, FitNet, AT, SP, CC, RKD) on the CIFAR-100 dataset under standardized conditions; and (3) practical observations on when symmetric vs.\ asymmetric distillation may be preferred.

%%%%%%%%%%%%%%%%%%%%%%%%%%%%%%%%%%%%%%%%%%
\section{Related Work}

\subsection{Knowledge Distillation}

Knowledge distillation~\cite{ref7} trains a student network to match the softened output distribution of a teacher network. The student loss combines standard cross-entropy with KL divergence:

\begin{equation}
L_{KD} = \alpha \cdot L_{CE}(y, \sigma(z_s)) + (1-\alpha) \cdot \tau^2 \cdot KL(\sigma(z_t/\tau) \parallel \sigma(z_s/\tau))
\end{equation}

where $z_s$ and $z_t$ are the student and teacher logits, $\tau$ is the temperature, and $\alpha$ balances the two terms.

Romero et al.~\cite{ref8} introduced FitNet, which adds an L2 loss between intermediate feature representations. Zagoruyko and Komodakis~\cite{ref9} proposed Attention Transfer (AT), which aligns spatial attention maps. These methods directly match specific representations between teacher and student.

A different line of research focuses on preserving relational structures. Park et al.~\cite{ref12} introduced Relational Knowledge Distillation (RKD), which preserves pairwise relationships. Tung and Mori~\cite{ref10} proposed Similarity Preserving (SP) distillation. Peng et al.~\cite{ref11} introduced Correlation Congruence (CC). These methods preserve structural relationships without requiring direct feature correspondence.

\subsection{Symmetry in Deep Learning}

The concept of symmetry has deep roots in deep learning. Convolutional neural networks exploit translational symmetry through weight sharing~\cite{ref13}. Group-equivariant CNNs~\cite{ref14} extend this to broader symmetry groups. Self-attention mechanisms~\cite{ref15} exhibit permutation symmetry. In representation learning, contrasting symmetric and asymmetric paradigms has proven fruitful: SimSiam~\cite{ref16} and BYOL~\cite{ref17} demonstrate that asymmetric architectures can prevent representational collapse. Our work draws a parallel distinction for knowledge distillation.

%%%%%%%%%%%%%%%%%%%%%%%%%%%%%%%%%%%%%%%%%%
\section{Symmetric vs.\ Asymmetric Knowledge Distillation}

\subsection{Formal Definition}

Let $f_T(x) \in \mathbb{R}^{d_T}$ and $f_S(x) \in \mathbb{R}^{d_S}$ denote the feature representations of the teacher and student networks. \textbf{Symmetric KD} minimizes a distance function $\mathcal{D}$ that directly compares representations on a per-sample, per-feature basis:

\begin{equation}
\mathcal{L}_{sym} = \mathbb{E}_{x \sim \mathcal{X}} \left[ \mathcal{D}\big(g_T(f_T(x)), g_S(f_S(x))\big) \right]
\end{equation}

\textbf{Asymmetric KD} minimizes a distance function that compares representations indirectly through their relational structures:

\begin{equation}
\mathcal{L}_{asym} = \mathbb{E}_{(x_i, x_j) \sim \mathcal{X}^2} \left[ \mathcal{D}\big(\mathcal{R}_T(f_T(x_i), f_T(x_j)), \mathcal{R}_S(f_S(x_i), f_S(x_j))\big) \right]
\end{equation}

\subsection{Symmetric KD Methods}

\textbf{KD (Hinton, 2015)} matches softened output distributions. \textbf{FitNet (Romero et al., 2015)} adds an L2 regression loss between intermediate features. \textbf{Attention Transfer (Zagoruyko \& Komodakis, 2017)} aligns normalized attention maps derived from feature activations.

\subsection{Asymmetric KD Methods}

\textbf{Similarity Preserving (Tung \& Mori, 2019)} preserves pairwise sample similarities. \textbf{Correlation Congruence (Peng et al., 2020)} aligns second-order feature correlations. \textbf{Relational Knowledge Distillation (Park et al., 2019)} preserves pairwise distance relationships.

%%%%%%%%%%%%%%%%%%%%%%%%%%%%%%%%%%%%%%%%%%
\section{Experimental Setup}

\subsection{Dataset and Preprocessing}

We conduct all experiments on the CIFAR-100 dataset~\cite{ref19}, consisting of 60,000 32$\times$32 color images across 100 classes. Standard data augmentation is applied: random crops with 4-pixel padding, random horizontal flips, and color jitter.

\subsection{Network Architectures}

\textbf{Teacher network:} ResNet-34 (21.3M parameters), trained from scratch for 60 epochs (best accuracy: 76.75\%). \textbf{Student network:} ResNet-18 (11.2M parameters), trained from random initialization.

\subsection{Training Protocol}

\begin{table}[H]
\caption{Training hyperparameters.\label{tab:hyperparams}}
\begin{tabularx}{\textwidth}{lC}
\toprule
\textbf{Hyperparameter} & \textbf{Value}\\
\midrule
Optimizer & SGD with Nesterov momentum\\
Momentum & 0.9\\
Weight decay & 5e-4\\
Batch size & 128\\
Epochs & 60\\
Initial learning rate & 0.05\\
LR schedule & Cosine annealing\\
Warmup epochs & 5\\
\bottomrule
\end{tabularx}
\end{table}

\begin{table}[H]
\caption{KD method hyperparameters.\label{tab:kdparams}}
\begin{tabularx}{\textwidth}{lCC}
\toprule
\textbf{Method} & \textbf{Category} & \textbf{Hyperparameters}\\
\midrule
KD~\cite{ref7} & Symmetric & $\tau=4.0$, $\alpha=0.5$\\
FitNet~\cite{ref8} & Symmetric & $\beta=1000$\\
AT~\cite{ref9} & Symmetric & $\beta=500$\\
SP~\cite{ref10} & Asymmetric & $\beta=3000$\\
CC~\cite{ref11} & Asymmetric & $\beta=0.02$, $\gamma=0.4$\\
RKD~\cite{ref12} & Asymmetric & $\beta=1.0$\\
\bottomrule
\end{tabularx}
\end{table}

\subsection{Evaluation Protocol}

All experiments are conducted on a single NVIDIA RTX 3090 GPU. Each method is trained for exactly 60 epochs with matched random seeds. We report the best test accuracy and full learning curves. \textbf{Limitation:} Results are based on single runs per method; multi-run experiments with statistical testing would strengthen the conclusions.

%%%%%%%%%%%%%%%%%%%%%%%%%%%%%%%%%%%%%%%%%%
\section{Results and Analysis}

\subsection{Main Results}

Table~\ref{tab:results} presents the best test accuracies achieved by each method.

\begin{table}[H]
\caption{Best test accuracy (\%) on CIFAR-100. Teacher: ResNet-34 (76.75\%). Student: ResNet-18.\label{tab:results}}
\begin{tabularx}{\textwidth}{lCCC}
\toprule
\textbf{Method} & \textbf{Category} & \textbf{Best Acc. (\%)} & $\bm{\Delta}$ \textbf{vs. Baseline}\\
\midrule
Baseline (No KD) & --- & 76.48 & ---\\
KD (Hinton, 2015) & Symmetric & \textbf{77.71} & +1.23\\
FitNet (Romero, 2015) & Symmetric & 76.06 & $-$0.42\\
AT (Zagoruyko, 2017) & Symmetric & 77.34 & +0.86\\
SP (Tung, 2019) & Asymmetric & \textbf{77.85} & +1.37\\
RKD (Park, 2019) & Asymmetric & 76.43 & $-$0.05\\
CC (Peng, 2020) & Asymmetric & 75.81 & $-$0.67\\
\bottomrule
\end{tabularx}
\end{table}

Key observations: (1) three of six methods (KD, AT, SP) improve upon the baseline; (2) the best method is SP (asymmetric) at 77.85\%; (3) symmetric methods show more consistent improvements, while asymmetric methods exhibit higher variance; (4) three methods surpass the teacher's accuracy (76.75\%).

\subsection{Learning Curves}

Figure~\ref{fig:curves} shows test accuracy during training. Symmetric methods converge faster in early epochs, while asymmetric methods catch up in mid-training.

\begin{figure}[H]
\centering
\includegraphics[width=12.0cm]{../figures/learning_curves.png}
\caption{Test accuracy vs.\ training epoch for all KD methods.\label{fig:curves}}
\end{figure}

\begin{figure}[H]
\centering
\includegraphics[width=12.0cm]{../figures/bar_comparison.png}
\caption{Bar chart comparing best test accuracy.\label{fig:bar}}
\end{figure}

\begin{figure}[H]
\centering
\includegraphics[width=10.0cm]{../figures/category_comparison.png}
\caption{Distribution of accuracies by category.\label{fig:violin}}
\end{figure}

\subsection{Category Comparison}

Symmetric average: 77.04\% (range: 76.06--77.71\%). Asymmetric average: 76.70\% (range: 75.81--77.85\%). Symmetric methods show a tighter range (1.65 pp vs.\ 2.04 pp).

%%%%%%%%%%%%%%%%%%%%%%%%%%%%%%%%%%%%%%%%%%
\section{Discussion}
\label{sec:discussion}

\subsection{The Symmetry Spectrum}

Our taxonomy represents a spectrum rather than a strict dichotomy. \textbf{Attention Transfer} computes attention maps through spatial aggregation (relational) but aligns them directly (symmetric). \textbf{FitNet} uses a learned regression function, introducing an asymmetric component. \textbf{KD with temperature scaling} transforms distributions before comparison, emphasizing relative relationships. We recommend specifying which level (representation transformation vs.\ comparison operation) is being referred to.

\subsection{Why Symmetric Methods Converge Faster}

Symmetric methods provide per-sample gradient signals that directly push each student feature toward its teacher counterpart, creating a dense gradient field. Asymmetric methods provide gradients through relational structures depending on pairwise interactions. Symmetric KD explicitly enforces equivariance between teacher and student feature maps, while asymmetric KD preserves invariance properties. Equivariance constraints provide more per-sample gradient information.

\subsection{Why Asymmetric Methods Can Achieve Higher Peak Performance}

SP's superior performance (77.85\%) can be attributed to the flexibility of relational knowledge transfer. By preserving similarity structures rather than forcing direct alignment, SP allows the student to develop representations naturally suited to its capacity. However, we caution against over-interpreting small accuracy differences---the gap between KD (77.71\%) and SP (77.85\%) is only 0.14 percentage points.

\subsection{Limitations}

(1) Single dataset (CIFAR-100) and architecture family (ResNet). (2) Single-run experiments per method. (3) Hyperparameters adopted from prior literature without per-method tuning. (4) Limited method coverage (six methods).

%%%%%%%%%%%%%%%%%%%%%%%%%%%%%%%%%%%%%%%%%%
\section{Conclusions}

We introduced a perspective for understanding knowledge distillation methods through the lens of symmetry, categorizing methods as symmetric (direct alignment) or asymmetric (relational preservation). Our experimental comparison on CIFAR-100 reveals that symmetric methods converge faster while asymmetric methods can achieve higher peak performance. The best-performing method, SP, achieves 77.85\%, surpassing both the baseline and teacher. We discuss boundary cases and acknowledge important limitations, providing a useful conceptual tool for practitioners.

%%%%%%%%%%%%%%%%%%%%%%%%%%%%%%%%%%%%%%%%%%
\vspace{6pt}

\authorcontributions{Conceptualization, F.L.; methodology, F.L.; software, F.L.; validation, F.L.; formal analysis, F.L.; investigation, F.L.; resources, F.L.; data curation, F.L.; writing---original draft preparation, F.L.; writing---review and editing, F.L.; visualization, F.L.; supervision, F.L.; project administration, F.L.; funding acquisition, F.L. All authors have read and agreed to the published version of the manuscript.}

\funding{This research received no external funding.}

\institutionalreview{Not applicable.}

\informedconsent{Not applicable.}

\dataavailability{The CIFAR-100 dataset is publicly available. The source code for all experiments will be made available upon request.}

\acknowledgments{The authors acknowledge the computational resources provided by NVIDIA RTX 3090 GPU.}

\conflictsofinterest{The authors declare no conflicts of interest.}

\abbreviations{Abbreviations}{%
\begin{tabular}{@{}ll}
KD & Knowledge Distillation\\
AT & Attention Transfer\\
SP & Similarity Preserving\\
CC & Correlation Congruence\\
RKD & Relational Knowledge Distillation\\
KL & Kullback--Leibler\\
CNN & Convolutional Neural Network\\
\end{tabular}
}

%%%%%%%%%%%%%%%%%%%%%%%%%%%%%%%%%%%%%%%%%%
\reftitle{References}

\begin{thebibliography}{20}

\bibitem{ref1}
He, K.; Zhang, X.; Ren, S.; Sun, J. Deep residual learning for image recognition. In \textit{Proceedings of the IEEE Conference on Computer Vision and Pattern Recognition}, 2016; pp. 770--778.

\bibitem{ref2}
Huang, G.; Liu, Z.; Van Der Maaten, L.; Weinberger, K.Q. Densely connected convolutional networks. In \textit{Proceedings of the IEEE Conference on Computer Vision and Pattern Recognition}, 2017; pp. 4700--4708.

\bibitem{ref3}
Tan, M.; Le, Q.V. EfficientNet: Rethinking model scaling for convolutional neural networks. In \textit{International Conference on Machine Learning}, PMLR, 2019; pp. 6105--6114.

\bibitem{ref4}
Cheng, Y.; Wang, D.; Zhou, P.; Zhang, T. A survey of model compression and acceleration for deep neural networks. \textit{arXiv preprint arXiv:1710.09282}, 2017.

\bibitem{ref5}
Han, S.; Pool, J.; Tran, J.; Dally, W. Learning both weights and connections for efficient neural network. In \textit{Advances in Neural Information Processing Systems}, 2015; pp. 1135--1143.

\bibitem{ref6}
Jacob, B.; Kligys, S.; Chen, B.; Zhu, M.; Tang, M.; Howard, A.; Adam, H.; Kalenichenko, D. Quantization and training of neural networks for efficient integer-arithmetic-only inference. In \textit{Proceedings of the IEEE Conference on Computer Vision and Pattern Recognition}, 2018; pp. 2704--2713.

\bibitem{ref7}
Hinton, G.; Vinyals, O.; Dean, J. Distilling the knowledge in a neural network. \textit{arXiv preprint arXiv:1503.02531}, 2015.

\bibitem{ref8}
Romero, A.; Ballas, N.; Kahou, S.E.; Chassang, A.; Gatta, C.; Bengio, Y. FitNets: Hints for thin deep nets. \textit{arXiv preprint arXiv:1412.6550}, 2014.

\bibitem{ref9}
Zagoruyko, S.; Komodakis, N. Paying more attention to attention: Improving the performance of convolutional neural networks via attention transfer. \textit{arXiv preprint arXiv:1612.03928}, 2016.

\bibitem{ref10}
Tung, F.; Mori, G. Similarity-preserving knowledge distillation. In \textit{Proceedings of the IEEE/CVF International Conference on Computer Vision}, 2019; pp. 1365--1374.

\bibitem{ref11}
Peng, B.; Jin, X.; Liu, J.; Li, D.; Wu, Y.; Liu, Y.; Zhou, S.; Zhang, Z. Correlation congruence for knowledge distillation. In \textit{Proceedings of the IEEE/CVF International Conference on Computer Vision}, 2019; pp. 5007--5016.

\bibitem{ref12}
Park, W.; Kim, D.; Lu, Y.; Cho, M. Relational knowledge distillation. In \textit{Proceedings of the IEEE/CVF Conference on Computer Vision and Pattern Recognition}, 2019; pp. 3967--3976.

\bibitem{ref13}
LeCun, Y.; Bottou, L.; Bengio, Y.; Haffner, P. Gradient-based learning applied to document recognition. \textit{Proceedings of the IEEE} \textbf{1998}, \textit{86}, 2278--2324.

\bibitem{ref14}
Cohen, T.; Welling, M. Group equivariant convolutional networks. In \textit{International Conference on Machine Learning}, PMLR, 2016; pp. 2990--2999.

\bibitem{ref15}
Vaswani, A.; Shazeer, N.; Parmar, N.; Uszkoreit, J.; Jones, L.; Gomez, A.N.; Kaiser, \L.; Polosukhin, I. Attention is all you need. In \textit{Advances in Neural Information Processing Systems}, 2017; pp. 5998--6008.

\bibitem{ref16}
Chen, X.; He, K. Exploring simple siamese representation learning. In \textit{Proceedings of the IEEE/CVF Conference on Computer Vision and Pattern Recognition}, 2021; pp. 15750--15758.

\bibitem{ref17}
Grill, J.B.; Strub, F.; Altch\'{e}, F.; Tallec, C.; Richemond, P.H.; Buchatskaya, E.; Doersch, C.; Avila Pires, B.; Guo, Z.D.; others. Bootstrap your own latent: A new approach to self-supervised learning. In \textit{Advances in Neural Information Processing Systems}, 2020; pp. 21271--21284.

\bibitem{ref18}
Tian, Y.; Krishnan, D.; Isola, P. Contrastive representation distillation. \textit{arXiv preprint arXiv:1910.10699}, 2019.

\bibitem{ref19}
Krizhevsky, A.; Hinton, G. Learning multiple layers of features from tiny images. \textit{Technical Report}, University of Toronto, 2009.

\bibitem{ref20}
Loshchilov, I.; Hutter, F. SGDR: Stochastic gradient descent with warm restarts. \textit{arXiv preprint arXiv:1608.03983}, 2016.

\end{thebibliography}

\end{document}
